\clearpage
\begin{explanationbox}
    \textbf{\textcolor{CExplanation}{一、为什么会出现这个结论？}}
    例如 $y=2x+1$，规定 $0\le x\le3$。题目问的是：$y$ 最小是多少、最大是多少、
    中间又能取到哪些数？如果每次都代入许多 $x$ 去试，既慢又不可能试完。
    注意到 $x$ 越大，$2x+1$ 也越大，于是直接计算
    \[
        f(0)=1,\qquad f(3)=7.
    \]
    这样就不用一个一个地试值了。一次函数图像是一条不断开的直线，
    从高度 $1$ 走到高度 $7$，中间的高度也都会经过，所以值域是 $[1,7]$。

    \medskip
    \textbf{\textcolor{CExplanation}{二、先分清两个不同的问题}}
    \begin{description}[style=nextline,leftmargin=2.8em,labelsep=0.5em,itemsep=0.45em]
        \item[\textbf{问题一（谁最大、谁最小？）：}]
            只要知道在闭区间上一直上升或一直下降，答案就藏在两个端点里。
            这一步只需要\textbf{单调}，不要求图像一定连起来。
        \item[\textbf{问题二（中间每个数都能取到吗？）：}]
            这一步还要检查\textbf{连续}。图像没有断开，才能确定最大值和最小值之间的高度都不会漏掉。
    \end{description}

    \medskip
    \textbf{\textcolor{CExplanation}{三、从图像上怎么理解？}}
    首页蓝色高度带显示连续递增时的完整值域；易错图的浅橙色带显示跳跃后缺少的高度。
“单调”确定两端的最值，“连续”保证两端之间不缺值。

    \medskip
    \textbf{\textcolor{CExplanation}{四、考场上为什么有用？}}
    \begin{itemize}[leftmargin=*,itemsep=0.35em,topsep=0.35em]
        \item \textbf{直接求值域：}一次函数、指定区间内的二次函数、指数函数等，确认单调后代入端点。
        \item \textbf{求最大值或最小值：}不用把整个函数解出来，只需比较两端。
        \item \textbf{有字母的题：}先看字母让函数上升还是下降，再决定端点排列顺序。
        \item \textbf{分段求范围：}二次函数跨过对称轴、绝对值函数遇到转折点时，先拆后合并。
    \end{itemize}

    \medskip
    \textbf{\textcolor{CExplanation}{五、两个容易让人眼前一亮的发现}}
    \begin{description}[style=nextline,leftmargin=2.8em,labelsep=0.5em,itemsep=0.45em]
        \item[\textbf{发现一（先比较大小，再写值域）：}]
            对连续且单调的函数，可以直接记
            \[
                \boxed{\text{值域}=
                [\min\{f(a),f(b)\},\ \max\{f(a),f(b)\}]}.
            \]
            这里 $\min$ 表示取较小的数，$\max$ 表示取较大的数。
            这样即使暂时忘了递增、递减的公式，也不容易写反。
        \item[\textbf{发现二（两个端点相等时要警惕）：}]
            如果整个区间上单调且 $f(a)=f(b)$，那么区间内的函数值只能都等于这个数。
            可是如果函数中途改变了方向，两端相等\textbf{不代表}中间也一样。
            例如 $f(x)=x^2$ 在 $[-1,1]$ 上两端都等于 $1$，最小值却是 $0$。
    \end{description}

    \medskip
    \textbf{\textcolor{CExplanation}{做题顺序}}
    \textbf{先看区间是否包含端点 $\longrightarrow$ 再看升降方向
    $\longrightarrow$ 算端点值 $\longrightarrow$ 检查有没有跳跃或转折 $\longrightarrow$ 写完整值域。}
\end{explanationbox}
